\documentclass[11pt]{article}

\usepackage[margin=1.1in]{geometry}
\usepackage{amsmath,amssymb,amsthm}
\usepackage{booktabs}
\usepackage{array}
\usepackage{microtype}
\usepackage{xcolor}
\usepackage[colorlinks=true,linkcolor=blue!60!black,citecolor=blue!60!black,urlcolor=blue!60!black]{hyperref}

\newtheorem{theorem}{Theorem}[section]
\newtheorem{lemma}[theorem]{Lemma}
\newtheorem{proposition}[theorem]{Proposition}
\newtheorem{corollary}[theorem]{Corollary}
\newtheorem{conjecture}[theorem]{Conjecture}
\theoremstyle{definition}
\newtheorem{definition}[theorem]{Definition}
\newtheorem{remark}[theorem]{Remark}
\newtheorem{question}[theorem]{Question}

\DeclareMathOperator{\rank}{rank}
\DeclareMathOperator{\fix}{fix}
\DeclareMathOperator{\Ker}{ker}
\newcommand{\ZZ}{\mathbb{Z}}
\newcommand{\QQ}{\mathbb{Q}}
\newcommand{\CC}{\mathbb{C}}
\newcommand{\NN}{\mathbb{N}}
\newcommand{\FF}{\mathbb{F}}
\newcommand{\seqA}{\textnormal{A000602}}
\newcommand{\seqT}{\textnormal{A000598}}

\title{Certified obstructions to a closed-form expression for the\\
number of alkane isomers:\\
\large computer-assisted non-existence certificates for the sequence A000602}
\author{}
\date{June 12, 2026}

\begin{document}
\maketitle

\begin{abstract}
Let $a(n)$ denote the number of structural isomers of the alkane
$\mathrm{C}_n\mathrm{H}_{2n+2}$, i.e.\ the number of unlabeled, unrooted trees on
$n$ vertices with maximum degree $4$ (OEIS \seqA). No exact closed form for
$a(n)$ is known. We localize the obstruction and erect certified, exact-arithmetic
perimeters around every mechanism that has historically produced exact formulas
for nontrivial counting sequences. Specifically, we prove by rigorous modular
rank certificates, validated against $621$ exactly computed terms (cross-checked
against the independently computed OEIS $b$-file), that: (i)~$a(n)$ satisfies
\emph{no} polynomial-coefficient recurrence, homogeneous or inhomogeneous, in a
staircase region extending to order~$2$/degree~$150$, order~$10$/degree~$46$,
and order~$304$ with constant coefficients --- consequently no fixed-depth proper
hypergeometric multisum closed form (Wilf--Zeilberger class) with canonical
annihilator in that region can equal $a(n)$; (ii)~the generating function
satisfies no algebraic equation of $y$-degree $\le 28$ within an analogous degree
staircase, and no first-order algebraic differential equation of degree $\le 7$
in each of $x$, $y$, $y'$; (iii)~$a(n)\bmod 2$ has $2$-kernel of size $\ge 31$
(maximal growth through depth $4$) and $a(n)\bmod m$ is not eventually periodic
with period $\le 200$ for $2\le m\le 7$, blocking digit-automatic closed forms.
We then prove that the \emph{labeled} analogue of the problem admits an exact
two-index finite-sum closed form, locating the obstruction precisely in the
unlabeled quotient: Burnside's lemma expresses $a(n)$ as a sum over the
partitions of $n$, an index set of unbounded dimension. Finally we show
structurally why the one known mechanism for collapsing such sums to finite
exact formulas --- Rademacher's modular circle method --- is unavailable: the
generating function is governed by a mixed $(2,3)$-Mahler-type functional system
whose singularities cascade toward a (conjecturally natural) boundary on the
unit circle without any known transformation theory. Together these results give
sharp necessary conditions that any future closed form must satisfy.
\end{abstract}

\tableofcontents

%======================================================================
\section{Introduction}
%======================================================================

\subsection{The problem}

For $n\ge 1$ let $a(n)$ denote the number of isomorphism classes of trees on $n$
unlabeled vertices in which every vertex has degree at most $4$. Chemically,
$a(n)$ is the number of constitutional (structural) isomers of the alkane
$\mathrm{C}_n\mathrm{H}_{2n+2}$, stereoisomers not distinguished. The sequence
begins (with the convention $a(0)=1$)
\[
1,\;1,\;1,\;1,\;2,\;3,\;5,\;9,\;18,\;35,\;75,\;159,\;355,\;802,\;1858,\;4347,\;\dots
\]
and is sequence \seqA{} of the OEIS \cite{OEIS}. Its enumeration, initiated by
Cayley in 1875 \cite{Cayley1875} and completed as an \emph{algorithm} by
P\'olya in 1937 \cite{Polya1937}, was the founding application of P\'olya
enumeration theory. Yet after a century and a half the following question is
open.

\begin{question}[Problem (A)]\label{q:A}
Does there exist a function $F(n)$, built from finitely many elementary
operations $(+,-,\times,\div)$, exponentials and logarithms, standard special
functions (e.g.\ $\Gamma$, $\zeta$, hypergeometric ${}_pF_q$), and \emph{finite}
(fixed-depth) sums and products --- directly evaluable at any $n$ without
knowledge of $a(1),\dots,a(n-1)$, and not an implicit functional equation or a
coefficient-extraction from a generating function --- such that $F(n)=a(n)$ for
all $n\ge 1$?
\end{question}

This note does not answer Question~\ref{q:A}. Its purpose is to demonstrate,
with certificates that are rigorous modulo only the correctness of widely
validated exact-arithmetic computations, that \emph{every known route to such a
formula is blocked}, to quantify exactly \emph{how far} each route is blocked,
and to isolate the unique structural obstruction. All computations are
reproducible from the scripts described in Appendix~\ref{app:repro}.

\subsection{Why the question is sharp}

Two remarks calibrate what ``closed form'' can possibly mean here.

\begin{proposition}[Additive approximation is equivalent to exactness]
\label{prop:rounding}
Suppose $F$ is directly evaluable and $|F(n)-a(n)|<\tfrac12$ for all $n\ge n_0$.
Then $\lfloor F(n)+\tfrac12\rfloor$ is an exact closed form for $a(n)$ for
$n \ge n_0$. Hence ``closed form with additive error $O(n^{-k})$'' is not a
relaxation of Question~\ref{q:A} but a reformulation of it.
\end{proposition}

\begin{proof}
Immediate from $a(n)\in\ZZ$.
\end{proof}

\begin{remark}
Question~\ref{q:A} is \emph{not} vacuous for sequences of this general type:
Rademacher's convergent series for the partition function $p(n)$
\cite{Rademacher1937}, truncated at $O(\sqrt n)$ terms with a certified error
below $\tfrac12$, is exactly such a directly evaluable formula, even though the
generating function $\prod (1-x^k)^{-1}$ has the unit circle as a natural
boundary. The mechanism enabling it --- modularity of the Dedekind eta function
--- will reappear in \S\ref{sec:structure} as precisely what $a(n)$ lacks.
\end{remark}

By contrast, the \emph{multiplicative} (relative-error) version of the question
is completely solved by classical singularity analysis: P\'olya \cite{Polya1937}
proved $a(n)\sim C\,\alpha^{n}n^{-5/2}$, and the expansion extends to every
order (see \S\ref{sec:structure}),
\begin{equation}\label{eq:asym}
a(n)\;=\;C\,\alpha^{n}\,n^{-5/2}\Bigl(1+\frac{c_1}{n}+\frac{c_2}{n^2}+\cdots+
\frac{c_m}{n^m}+O(n^{-m-1})\Bigr),
\end{equation}
with $\alpha=\rho^{-1}=2.815460033176150746526616778\ldots$,
$C=0.656318695841834750\ldots$, $c_1=0.03006869504\ldots$,
$c_2=4.7882749807\ldots$, $c_3\approx 22.18$ (constants computed in this
project; $\alpha$ and $C$ classical \cite{RHB1976,Finch2003,FS2009}). The gap
between \eqref{eq:asym} and Question~\ref{q:A} is exponential: by
Proposition~\ref{prop:rounding} an exact formula needs relative accuracy
$\alpha^{-n}n^{5/2}$, while any fixed truncation of \eqref{eq:asym} achieves
only polynomial relative accuracy.

\subsection{Summary of results}

Write $\mathcal{R}(r,d)$ for the statement ``there exist polynomials
$p_0,\dots,p_r,q\in\CC[n]$ of degree $\le d$, with $(p_0,\dots,p_r)\neq 0$, such
that $\sum_{j=0}^{r}p_j(n)\,a(n+j)=q(n)$ for all $n\ge 1$''.

\begin{theorem}[Non-holonomicity perimeter; \S\ref{sec:holo}]\label{thm:main-holo}
$\mathcal{R}(r,d)$ is \emph{false} for every pair $(r,d)$ in the staircase
region of Table~\ref{tab:holo}, which contains in particular
$(2,150)$, $(3,120)$, $(4,100)$, $(10,46)$, $(56,8)$, $(110,3)$, $(200,1)$ and
$(304,0)$.
\end{theorem}

\begin{theorem}[Non-algebraicity perimeter; \S\ref{sec:alg}]\label{thm:main-alg}
The generating function $A(x)=\sum_{n\ge0}a(n)x^n$ satisfies no nonzero
polynomial relation $P(x,A(x))=0$ with $\deg_y P\le D_y$, $\deg_x P\le D_x$ for
any $(D_y,D_x)$ in the staircase of Table~\ref{tab:alg} (up to $(28,19)$ and
$(2,190)$), and no relation $P(x,A(x),A'(x))=0$ with degree $\le 7$ in each
variable.
\end{theorem}

\begin{theorem}[Arithmetic obstructions; \S\ref{sec:autom}]\label{thm:main-autom}
The $2$-kernel of $(a(n)\bmod 2)$ has at least $31$ pairwise distinct elements
(its growth through depth $4$ is the maximum possible, $2^{k+1}-1$), and the
$3$-kernel of $(a(n)\bmod 3)$ has at least $40$. For every modulus $2\le m\le 7$,
$(a(n)\bmod m)$ is not eventually periodic with any period $\le 200$.
\end{theorem}

\begin{theorem}[The labeled problem is closed-form; \S\ref{sec:labeled}]
\label{thm:main-labeled}
The number $L(n)$ of \emph{labeled} trees on $n\ge3$ vertices with maximum
degree $4$ equals the two-index finite sum
\[
L(n)=\sum_{d=0}^{\lfloor (n-2)/3\rfloor}\;
\sum_{c=0}^{\lfloor (n-2-3d)/2\rfloor}
\frac{n!\,(n-2)!}{a!\;b!\;c!\;d!\;2^{c}\,6^{d}},
\qquad b=n-2-2c-3d,\quad a=n-b-c-d ,
\]
which lies squarely in the function class of Question~\ref{q:A}. Hence the
degree-$4$ (valence) constraint is \emph{not} the obstruction; the unlabeled
quotient is.
\end{theorem}

Section~\ref{sec:structure} assembles the structural picture: the Burnside sum
over partitions, the Otter cancellation that produces the exponent $-5/2$, the
singularity cascade of the underlying mixed $(2,3)$-Mahler system, and the
absence of a Rademacher mechanism. Section~\ref{sec:necessary} states the
resulting necessary conditions on any solution of Question~\ref{q:A}.

\subsection*{Epistemic status}
Theorems~\ref{thm:main-holo}--\ref{thm:main-autom} are computer-assisted. The
underlying linear algebra is exact (machine-word arithmetic modulo primes
$<2^{26}$, no floating point); the logical step from a full-rank computation to
non-existence is a one-line lemma (Lemma~\ref{lem:rank}); the data feeding the
matrices is certified in \S\ref{sec:data}. Implementation risk is controlled by
decoy sequences with known positive answers (the guessers \emph{do} find the
Catalan recurrence, the Catalan algebraic equation, and the Thue--Morse kernel)
and by re-verification modulo a second prime. The propositions of
\S\ref{sec:structure} marked ``sketch'' are structural/heuristic and clearly
flagged; nothing in \S\S\ref{sec:holo}--\ref{sec:labeled} depends on them.

%======================================================================
\section{The defining system and certified data}\label{sec:data}
%======================================================================

\subsection{The P\'olya--Otter system}

Let $Z_{S_3}$ and $Z_{S_4}$ be the cycle index polynomials
\[
Z_{S_3}(t_1,t_2,t_3)=\tfrac{1}{6}\bigl(t_1^3+3t_1t_2+2t_3\bigr),
\qquad
Z_{S_4}(t_1,t_2,t_3,t_4)=\tfrac{1}{24}\bigl(t_1^4+6t_1^2t_2+3t_2^2+8t_1t_3+6t_4\bigr).
\]
Let $T(x)=\sum_{n\ge0}t(n)x^n$ be the generating function of \seqT{} (alkyl
radicals: rooted trees in which every vertex has at most $3$ children), defined
by
\begin{equation}\label{eq:T}
T(x)\;=\;1+x\,Z_{S_3}\bigl(T(x),T(x^2),T(x^3)\bigr)
\;=\;1+\frac{x}{6}\Bigl(T(x)^3+3\,T(x)\,T(x^2)+2\,T(x^3)\Bigr).
\end{equation}
The unrooted count is recovered by the dissymmetry (Otter) formula
\cite{Otter1948,HararyNorman1960,RainsSloane1999}:
\begin{equation}\label{eq:A}
A(x)\;=\;\sum_{n\ge0}a(n)x^n
\;=\;1+x\,Z_{S_4}\bigl(T(x),T(x^2),T(x^3),T(x^4)\bigr)
-\frac{\bigl(T(x)-1\bigr)^2-\bigl(T(x^2)-1\bigr)}{2}.
\end{equation}
Equations \eqref{eq:T}--\eqref{eq:A} are the standard formulation (they appear,
in equivalent form, on the OEIS pages of \seqT{} and \seqA{}); they constitute a
\emph{recursive} description, not a closed form: extracting $a(n)$ requires the
$n$ preceding coefficients of $T$.

\subsection{Exact computation and certification}\label{subsec:cert}

Working in exact integer arithmetic, the coefficients $t(0),\dots,t(620)$ were
computed from \eqref{eq:T} by incremental convolution (the divisions by $6$ are
exact at every step, a useful integrality check that was asserted at each $n$),
and $a(0),\dots,a(620)$ from \eqref{eq:A}. The value $a(620)$ has $272$ decimal
digits.

\medskip
\noindent\textbf{Certification of the data.}
\begin{enumerate}
\item $t(0..14)$ and $a(0..20)$ agree with the values printed in the OEIS
entries \seqT{}, \seqA{}.
\item All $621$ computed values $a(0),\dots,a(620)$ agree with the OEIS
$b$-file \texttt{b000602.txt} (terms $n=0..1000$, contributed independently by
P.~von Br\"omssen), downloaded and compared term-by-term.
\item As an end-to-end validation of the same code path on a published
high-precision target: the singularity-analysis pipeline built on the computed
$t(n)$ reproduces Kotesovec's $30$-digit constant for \seqT{},
$t(n)\sim C_R\,\alpha^{n}n^{-3/2}$ with
$C_R=0.51787590645889353699\ldots$, to $22$ significant digits.
\end{enumerate}
All matrices in \S\S\ref{sec:holo}--\ref{sec:alg} are populated from these
certified integers reduced modulo the primes
\begin{equation}\label{eq:primes}
p_1=67108859=2^{26}-5,\qquad p_2=67108837 ,
\end{equation}
chosen below $2^{26}$ so that all modular row operations stay far inside the
exact range of $64$-bit integer arithmetic.

%======================================================================
\section{The certificate method}\label{sec:method}
%======================================================================

All non-existence theorems below reduce to the statement that an explicit
integer matrix has \emph{full column rank}. The logic is collected here once.

\begin{lemma}[Modular rank certificate]\label{lem:rank}
Let $M\in\ZZ^{R\times S}$ and let $p$ be prime. If the reduction
$\overline{M}\in\FF_p^{R\times S}$ has $\rank_{\FF_p}\overline{M}=S$, then
$\rank_\QQ M=S$; consequently $Mv=0$ has no nonzero solution $v\in\CC^{S}$.
\end{lemma}

\begin{proof}
Reduction mod $p$ cannot increase rank: a nonvanishing $S\times S$ minor of
$\overline M$ lifts to a minor of $M$ not divisible by $p$, hence nonzero. So
$\rank_\QQ M = S$ and $\Ker_\QQ M=0$. Since $M$ has rational entries,
$\dim_\CC\Ker_\CC M=\dim_\QQ\Ker_\QQ M=0$.
\end{proof}

\begin{remark}\label{rem:onesided}
The certificate is one-sided in exactly the safe direction: a \emph{rank
deficiency} mod $p$ could be an accident of the prime, but \emph{full rank} mod a
single prime is a rigorous proof of non-existence over $\CC$. The second prime
$p_2$ in \eqref{eq:primes} guards only against implementation error, not against
a gap in the mathematics. Note also that Lemma~\ref{lem:rank} yields
non-existence over $\CC$, not merely over $\QQ$: closed forms with irrational or
transcendental constants ($\pi$, $\zeta(3)$, algebraic irrationals, \dots) are
covered automatically.
\end{remark}

\begin{lemma}[Window reduction: eventual relations imply windowed relations]
\label{lem:window}
Suppose $\sum_{j=0}^{r}p_j(n)a(n+j)=q(n)$ holds for all $n\ge n_0$, with
$\deg p_j,\deg q\le d$ and $(p_j)\neq0$. Put
$s(n)=\prod_{i=1}^{n_0-1}(n-i)$. Then
$\sum_{j=0}^{r}\bigl(s\,p_j\bigr)(n)\,a(n+j)=\bigl(s\,q\bigr)(n)$ holds for
\emph{all} $n\ge1$, with degrees $\le d+n_0-1$. Hence ruling out relations valid
from $n=1$ in a degree-budgeted region also rules out eventually-valid relations
with correspondingly smaller degree.
\end{lemma}

\begin{proof}
For $n\ge n_0$ multiply the hypothesis by $s(n)$; for $1\le n<n_0$ both sides
vanish because $s(n)=0$.
\end{proof}

\begin{lemma}[Finite windows suffice]\label{lem:finitewindow}
Fix $(r,d)$ and let $S=(r+1)(d+1)+(d+1)$ count the unknown coefficients of
$(p_0,\dots,p_r,q)$. Build $M\in\ZZ^{R\times S}$ whose row indexed by $n$
($1\le n\le R$) expresses $\sum_j p_j(n)a(n+j)-q(n)=0$ in those unknowns. If
$\rank M=S$ for some finite $R$, then $\mathcal{R}(r,d)$ is false: no relation
holds for all $n\ge1$, since it would already have to hold on the window.
\end{lemma}

\begin{proof}
A relation valid for all $n\ge1$ is in particular valid for $1\le n\le R$, i.e.\
its coefficient vector lies in $\Ker M=0$.
\end{proof}

In all runs the window length was $R=\min(620-r,\;S+25)$, which always exceeded
$S$ by at least $8$ rows; ranks were computed by exact Gaussian elimination over
$\FF_{p_1}$ (and re-verified over $\FF_{p_2}$ for three representative cells).

\medskip
\noindent\textbf{Decoy validation.} Before each search the identical code was
required to \emph{find} a known relation: the Catalan numbers
($(n+2)C_{n+1}=(4n+2)C_n$; the guesser reports rank deficiency exactly $1$ at
$(r,d)=(1,1)$), the Catalan generating function ($y=1+xy^2$; found at
$\deg_y=2$), and the Thue--Morse sequence (whose $2$-kernel closes with $2$
states). A search rig that finds the planted relations and reports full rank on
$a(n)$ is making a meaningful assertion.

%======================================================================
\section{No polynomial recurrence: the holonomic perimeter}\label{sec:holo}
%======================================================================

\subsection{Why this class matters: the Wilf--Zeilberger reduction}

\begin{definition}[Hypergeometric-sum closed forms]\label{def:class}
Let $\mathcal{C}$ be the class of expressions $F(n)$ built, with fixed structure
independent of $n$, from:
(i) rational functions of $n$ with constants in $\CC$ (the constants may involve
$\pi$, $\zeta(s)$ at fixed $s$, $\Gamma$ at fixed arguments, etc.);
(ii) finitely many factors $\alpha^{\,n}$ ($\alpha\in\CC$ fixed) and
$\Gamma(\lambda n+\mu)^{\pm1}$, binomials $\binom{\lambda n+\mu}{\nu n+\xi}$,
with $\lambda,\nu\in\QQ$, $\mu,\xi\in\CC$;
(iii) nested finite sums $\sum_{k=0}^{\lambda n+\mu}$, of fixed depth, of
\emph{proper hypergeometric terms} \cite{WilfZeilberger1992,PWZ1996} in $n$ and
the summation indices, including evaluations ${}_pF_q(\dots;z_0)$ at fixed
$z_0$ with parameters affine in $n$;
(iv) interleavings by residue classes of $n$ modulo a fixed integer.
\end{definition}

Class $\mathcal{C}$ contains the overwhelming majority of expressions matching
the operation list of Question~\ref{q:A}: every binomial-sum formula, every
$\Gamma$-quotient formula, every terminating-hypergeometric formula in the
literature of exact enumeration is in $\mathcal{C}$. (What escapes
$\mathcal{C}$: constructs with $n$ inside a \emph{non-affine} or non-arithmetic
position, e.g.\ $\lfloor n\theta\rfloor$ with $\theta$ irrational, $n^n$ inside a
summand bound, or $n$-dependent analytic parameters; see
\S\ref{sec:necessary}.)

\begin{theorem}[Wilf--Zeilberger; Lipshitz closure]\label{thm:WZ}
Every $F\in\mathcal{C}$ is P-recursive over $\CC$: there exist
$p_0,\dots,p_r\in\CC[n]$, not all zero, and $q\in\CC[n]$ (the inhomogeneous part
absorbing degenerate boundary contributions) with
$\sum_{j=0}^{r}p_j(n)F(n+j)=q(n)$ for all sufficiently large $n$.
\end{theorem}

\begin{proof}[Proof references]
Proper hypergeometric multisums are holonomic
\cite{WilfZeilberger1992,PWZ1996}; P-recursiveness is preserved under sums,
products, Cauchy products, and interleaving \cite{Stanley1980,Lipshitz1989};
$\Gamma$-factors with rational $\lambda$ are handled by passing to arithmetic
subprogressions and interleaving.
\end{proof}

Thus a single family of certificates --- non-existence of P-recurrences ---
simultaneously eliminates the entire hypergeometric core of the function class
of Question~\ref{q:A}.

\subsection{The certificates}

For each pair $(r,d)$, the matrix of Lemma~\ref{lem:finitewindow} (including the
inhomogeneous column block for $q$) was assembled from the certified values
$a(1),\dots,a(620)$ and its rank computed over $\FF_{p_1}$.

\begin{table}[ht]
\centering\small
\begin{tabular}{rrrrl@{\qquad\qquad}rrrrl}
\toprule
$r$ & $d$ & $S$ & $\rank$ & verdict & $r$ & $d$ & $S$ & $\rank$ & verdict\\
\midrule
  2 & 150 & 604 & 604 & none &  26 & 18 & 532 & 532 & none\\
  3 & 120 & 605 & 605 & none &  31 & 15 & 528 & 528 & none\\
  4 & 100 & 606 & 606 & none &  38 & 12 & 520 & 520 & none\\
  5 &  85 & 602 & 602 & none &  46 & 10 & 528 & 528 & none\\
  6 &  72 & 584 & 584 & none &  56 &  8 & 522 & 522 & none\\
  8 &  55 & 560 & 560 & none &  70 &  6 & 504 & 504 & none\\
 10 &  46 & 564 & 564 & none &  84 &  5 & 516 & 516 & none\\
 12 &  38 & 546 & 546 & none &  96 &  4 & 490 & 490 & none\\
 15 &  31 & 544 & 544 & none & 110 &  3 & 448 & 448 & none\\
 18 &  26 & 540 & 540 & none & 145 &  2 & 441 & 441 & none\\
 22 &  21 & 528 & 528 & none & 200 &  1 & 404 & 404 & none\\
    &     &     &     &      & 304 &  0 & 306 & 306 & none\\
\bottomrule
\end{tabular}
\caption{Non-existence certificates for inhomogeneous P-recurrences
$\sum_{j=0}^{r}p_j(n)a(n+j)=q(n)$, $\deg p_j,\deg q\le d$. $S$ is the number of
unknown coefficients; ``$\rank$'' is the exact rank over $\FF_{p_1}$ of the
window system. Full rank ($\rank=S$) proves non-existence over $\CC$
(Lemma~\ref{lem:rank}). Cells $(2,150)$, $(10,46)$, $(304,0)$ re-verified over
$\FF_{p_2}$.}\label{tab:holo}
\end{table}

\begin{theorem}\label{thm:holo}
For every $(r,d)$ dominated by an entry of Table~\ref{tab:holo} (i.e.\ $r\le r_i$
and $d\le d_i$ for some row $i$), there is no identity
$\sum_{j=0}^{r}p_j(n)\,a(n+j)=q(n)$ valid for all $n\ge1$ with
$p_j,q\in\CC[n]$ of degree $\le d$ and $(p_j)\neq0$. By Lemma~\ref{lem:window},
there is also no such identity valid for $n\ge n_0$ whenever $d+n_0-1$ stays
within the same region.
\end{theorem}

\begin{corollary}\label{cor:noWZ}
No $F\in\mathcal{C}$ whose minimal annihilating operator (in the inhomogeneous
sense of Theorem~\ref{thm:WZ}) has order and degree inside the region of
Theorem~\ref{thm:holo} can satisfy $F(n)=a(n)$ for all large $n$. In particular
$a(n)$ is given by no binomial/$\Gamma$/hypergeometric finite-sum formula of
order $\le 2$ and degree $\le 150$, order $\le 10$ and degree $\le 46$, \dots,
or constant-coefficient recurrence of order $\le 304$ (the last excluding, e.g.,
every ``exponential polynomial'' $\sum_i P_i(n)\beta_i^n$ with $\le 305$ terms).
\end{corollary}

\begin{remark}[Calibration of the region]
The certified staircase $(r+1)(d+1)\lesssim 600$ is enormous relative to
naturally occurring formulas: the minimal recurrences of virtually all classical
closed-form sequences (central binomials, Catalan, Motzkin, Ap\'ery, Domb,
families of lattice-path and constant-term sequences) have $(r+1)(d+1)\le 40$.
A hypothetical hypergeometric formula for $a(n)$ would need to be two orders of
magnitude more complex than any closed form ever encountered in enumerative
combinatorics --- while reproducing a sequence whose defining structure
(\S\ref{sec:structure}) predicts no recurrence at all.
\end{remark}

%======================================================================
\section{No algebraic equation, no first-order ADE}\label{sec:alg}
%======================================================================

The second classical reservoir of closed forms is Lagrange inversion: if
$A(x)$ were algebraic, $a(n)$ would inherit (by Comtet's theorem
\cite{Comtet1964} and the theory of diagonals) precisely the kind of
$\Gamma$-quotient asymptotics and frequently a finite-sum coefficient formula
(Catalan, Motzkin, Schr\"oder, \dots). We certify this away directly at the
level of $A$ itself.

For a candidate relation $P(x,y)=\sum_{i\le D_x,\,j\le D_y}c_{ij}x^iy^j$ with
$P(x,A(x))=0$, impose the vanishing of the coefficients of
$x^0,\dots,x^{620}$ in $P(x,A(x))$: a linear system in the
$S=(D_x+1)(D_y+1)$ unknowns $c_{ij}$, with integer entries computed from the
certified series ($y^j$ computed by exact modular convolution). Full rank again
proves non-existence over $\CC$.

\begin{table}[ht]
\centering\small
\begin{tabular}{rrrrl@{\qquad\qquad}rrrrl}
\toprule
$D_y$ & $D_x$ & $S$ & $\rank$ & verdict & $D_y$ & $D_x$ & $S$ & $\rank$ & verdict\\
\midrule
 2 & 190 & 573 & 573 & none & 12 & 44 & 585 & 585 & none\\
 3 & 145 & 584 & 584 & none & 14 & 38 & 585 & 585 & none\\
 4 & 115 & 580 & 580 & none & 17 & 32 & 594 & 594 & none\\
 5 &  95 & 576 & 576 & none & 20 & 27 & 588 & 588 & none\\
 6 &  82 & 581 & 581 & none & 24 & 22 & 575 & 575 & none\\
 8 &  64 & 585 & 585 & none & 28 & 19 & 580 & 580 & none\\
10 &  52 & 583 & 583 & none &    &    &     &     &     \\
\bottomrule
\end{tabular}
\caption{Non-existence certificates for algebraic relations $P(x,A(x))=0$,
$\deg_y\le D_y$, $\deg_x\le D_x$. The decoy ($y=1+xy^2$ for Catalan) is found by
the same code.}\label{tab:alg}
\end{table}

\begin{theorem}\label{thm:alg}
$A(x)$ satisfies no nonzero $P(x,A)=0$ with $(\deg_y,\deg_x)$ dominated by an
entry of Table~\ref{tab:alg}. Moreover, with
$P(x,y,z)=\sum_{i,j,k\le 7}c_{ijk}x^iy^jz^k$ ($512$ unknowns, window of $620$
coefficient equations), the system for $P(x,A(x),A'(x))=0$ has full rank $512$:
$A$ satisfies no first-order algebraic differential equation of degree $\le7$ in
each variable.
\end{theorem}

\begin{remark}
Theorem~\ref{thm:holo} already implies (within its region) that $A$ is not
D-finite with small annihilator, which subsumes small algebraic equations
\cite{Comtet1964}; Theorem~\ref{thm:alg} is nevertheless logically independent
and pushes the $y$-degree much further than the D-finite staircase reaches.
The first-order ADE exclusion targets closed forms of Abel/Lambert type
($y'=R(x,y)$) that are not D-finite.
\end{remark}

%======================================================================
\section{No digit formula: automaticity and congruences}\label{sec:autom}
%======================================================================

A third mechanism producing exact, directly evaluable formulas is
\emph{automaticity}/regularity: if $(a(n)\bmod m)$ were $k$-automatic, or $a(n)$
itself $k$-regular, then $a(n)$ (or its residues) would be computable by a fixed
finite-state device reading the base-$k$ digits of $n$ --- the Stern--Brocot,
Rudin--Shapiro, paperfolding paradigm \cite{AlloucheShallit2003}. The functional
system \eqref{eq:T}--\eqref{eq:A} has visible Mahler-type structure (the
substitutions $x\mapsto x^2,x^3,x^4$), so this possibility deserves a genuine
test rather than dismissal.

Recall the $k$-kernel of a sequence $s$: the set of subsequences
$n\mapsto s(k^e n+r)$, $0\le r<k^e$, $e\ge0$. $s$ is $k$-automatic iff its
$k$-kernel is finite \cite[Thm.~6.6.2]{AlloucheShallit2003}.

\begin{theorem}\label{thm:autom}
Within the certified data range, with subsequences compared on overlaps of
length $\ge 20$:
\begin{enumerate}
\item the $2$-kernel of $(a(n)\bmod 2)$ contains, through depth $4$,
$1,3,7,15,31$ pairwise distinct elements --- the theoretical maximum
$2^{e+1}-1$ at every depth $e\le4$; no two of the $31$ subsequences agree;
\item the $3$-kernel of $(a(n)\bmod 3)$ grows $1,4,13,40$ through depth $3$
(again maximal, $(3^{e+1}-1)/2$); likewise the $3$-kernel of $(a(n)\bmod2)$;
\item for each $m\in\{2,3,4,5,6,7\}$, no eventual period $P\le200$ exists for
$(a(n)\bmod m)$ on $250\le n\le 620$.
\end{enumerate}
Consequently any $2$-DFAO computing $a(n)\bmod 2$ needs $\ge31$ kernel states,
any $3$-DFAO for $a(n)\bmod3$ needs $\ge40$; the same Thue--Morse-validated code
confirms maximal-rate kernel growth, the canonical signature of
non-automaticity.
\end{theorem}

\begin{remark}[Honest scope]
Items (1)--(3) are finite certificates: they prove kernel \emph{lower bounds}
and rule out small periods; unbounded kernel growth (full non-automaticity) is
not finitely certifiable and is here an extrapolation --- but one supported by
maximal growth at every tested depth, in both bases, for both moduli. Note the
contrast with genuinely automatic decoys, whose kernels close at depth $1$--$2$.
By Cobham's theorem \cite{Cobham1969} simultaneous $2$- and $3$-automaticity
would force eventual periodicity, which (3) excludes up to period $200$.
\end{remark}

This blocks the last standard route to an exact formula for the \emph{residues}
of $a(n)$, let alone $a(n)$ itself; it also rules out the possibility that
$a(n)$ is one of the rare ``Mahler-automatic'' coincidences in which a
$(2,3)$-mixed system degenerates to a digit rule. That degeneration, by the
rigidity theory of Mahler equations (Adamczewski--Bell \cite{AdamczewskiBell2017},
Sch\"afke--Singer \cite{SchafkeSinger2019}), happens essentially only for
rational functions --- and $A(x)$ is certified non-rational with vast margin by
Theorem~\ref{thm:alg}.

%======================================================================
\section{The labeled problem has a closed form}\label{sec:labeled}
%======================================================================

The next theorem shows that the chemically meaningful constraint --- maximum
degree $4$ --- is \emph{not} what blocks Question~\ref{q:A}.

\begin{theorem}\label{thm:labeled}
For $n\ge3$, the number of labeled trees on $\{1,\dots,n\}$ with all degrees
$\le4$ is
\[
L(n)\;=\;\sum_{d=0}^{\lfloor (n-2)/3\rfloor}\;\;
\sum_{c=0}^{\lfloor (n-2-3d)/2\rfloor}
\frac{n!}{a!\,b!\,c!\,d!}\cdot\frac{(n-2)!}{2^{c}\,6^{d}},
\qquad
\begin{aligned}
b&=n-2-2c-3d,\\
a&=n-b-c-d,
\end{aligned}
\]
a fixed-depth (two-index) finite sum of $\Gamma$-quotients, i.e.\ a member of
class $\mathcal{C}$ and of the function class of Question~\ref{q:A}; moreover
$L\in\mathcal{C}$ is P-recursive by Theorem~\ref{thm:WZ}.
\end{theorem}

\begin{proof}
By the Pr\"ufer correspondence, vertex $i$ appears exactly $\deg(i)-1$ times in
the Pr\"ufer word, so the number of labeled trees with prescribed degrees
$d_1,\dots,d_n$ is the multinomial $\binom{n-2}{d_1-1,\dots,d_n-1}$
\cite{Moon1970}. Group the vertices by degree: $a$ vertices of degree $1$, $b$
of degree $2$, $c$ of degree $3$, $d$ of degree $4$. The handshake identity
$\sum d_i=2(n-1)$ together with $a+b+c+d=n$ forces $b+2c+3d=n-2$, leaving the
two free indices $c,d$. Choosing which labels receive which degree contributes
$n!/(a!\,b!\,c!\,d!)$; the Pr\"ufer multinomial contributes
$(n-2)!/\bigl(0!^{\,a}1!^{\,b}2!^{\,c}3!^{\,d}\bigr)=(n-2)!/(2^c6^d)$. Summing
over admissible $(c,d)$ gives the formula. It was verified against exhaustive
Pr\"ufer-word enumeration for $n\le8$:
$L(1..8)=1,1,3,16,125,1290,16590,255920$.
\end{proof}

\begin{remark}
The same argument gives a closed form for \emph{any} fixed degree bound, and
analogous exact formulas exist for labeled forests, labeled unicyclic graphs,
etc. The entire difficulty of Question~\ref{q:A} therefore resides in passing
from $L(n)$ to its unlabeled quotient.
\end{remark}

%======================================================================
\section{Localizing the obstruction}\label{sec:structure}
%======================================================================

\subsection{Burnside over partitions: unbounded index dimension}

The unlabeled count is the orbit count of the $S_n$-action on labeled trees:
\begin{equation}\label{eq:burnside}
a(n)\;=\;\frac{1}{n!}\sum_{\sigma\in S_n}\fix(\sigma)
\;=\;\sum_{\lambda\,\vdash\,n}\frac{\fix(\lambda)}{z_\lambda},
\end{equation}
where $\fix(\lambda)$ counts degree-$\le4$ labeled trees invariant under a
permutation of cycle type $\lambda$ and $z_\lambda$ is the usual centralizer
order. Equation \eqref{eq:burnside} \emph{is} an exact formula --- but its index
set is the set of partitions of $n$: $p(n)\sim e^{\pi\sqrt{2n/3}}/(4n\sqrt3)$
terms, indexed by integer vectors of \emph{unbounded length}. A fixed-depth
finite sum in the sense of Question~\ref{q:A} has an index set of bounded
dimension growing polynomially. The combinatorial content of
Theorems~\ref{thm:holo}--\ref{thm:autom} is that no collapse of
\eqref{eq:burnside} to bounded depth exists within the holonomic, algebraic, or
automatic worlds. (For trees, $\fix(\lambda)$ is supported on special $\lambda$
and is computed in practice through \eqref{eq:T}--\eqref{eq:A}; the partition
sum is a faithful image of the substitutions $x\mapsto x^m$ in
\eqref{eq:T}--\eqref{eq:A}.)

The one known precedent for such a collapse is instructive.

\begin{remark}[The Rademacher precedent]\label{rem:rademacher}
$p(n)$ itself is the orbit-count-like coefficient of a natural-boundary
generating function, yet possesses the exact formula
\[
p(n)=\frac{1}{\pi\sqrt2}\sum_{k=1}^{O(\sqrt n)}A_k(n)\,\sqrt{k}\;
\frac{d}{dn}\!\left(\frac{\sinh\bigl(\tfrac{\pi}{k}
\sqrt{\tfrac23(n-\tfrac1{24})}\bigr)}{\sqrt{n-\tfrac1{24}}}\right)
+\;\Theta,\qquad |\Theta|<\tfrac12,
\]
directly evaluable and exact after rounding \cite{Rademacher1937}. Every step of
its proof consumes the $\mathrm{SL}_2(\ZZ)$-modularity of
$x^{1/24}/\eta(\tau)$: the circle method needs an exact transformation law at
\emph{every} rational point of the boundary. Thus ``natural boundary''
per se is not the obstruction to Question~\ref{q:A}; \emph{natural boundary
without a transformation group} is.
\end{remark}

\subsection{The singular structure of the alkane system}

\subsubsection{The Otter cancellation and the exponent $-5/2$}

\begin{lemma}[Cycle-index derivative identity]\label{lem:cancel}
$\displaystyle \frac{\partial Z_{S_4}}{\partial t_1}=Z_{S_3}$.
\end{lemma}

\begin{proof}
$\partial_{t_1}\tfrac1{24}(t_1^4+6t_1^2t_2+3t_2^2+8t_1t_3+6t_4)
=\tfrac1{24}(4t_1^3+12t_1t_2+8t_3)=\tfrac16(t_1^3+3t_1t_2+2t_3)$.
(Equivalently: differentiating the cycle index of $S_4$ in the
fixed-points slot deletes a fixed point, leaving $S_3$.)
\end{proof}

\begin{proposition}[Square-root singularity and its cancellation]
\label{prop:otter}
$T$ has radius of convergence
$\rho=0.3551817423143773928822444736476\ldots$, with
$a:=T(\rho)=2.1174207009536310225162585905\ldots$ determined by the smooth
implicit system
\[
a=1+\rho\,Z_{S_3}\bigl(a,T(\rho^2),T(\rho^3)\bigr),
\qquad
\tfrac{\rho}{2}\bigl(a^2+T(\rho^2)\bigr)=1 ,
\]
and a local expansion
$T(x)=a-b\sqrt{1-x/\rho}+c\,(1-x/\rho)+d\,(1-x/\rho)^{3/2}+\cdots$ with $b>0$.
Writing \eqref{eq:A} as $A(x)=G\bigl(x,T(x)\bigr)$ with
$G(x,t)=1+x\,Z_{S_4}(t,T(x^2),T(x^3),T(x^4))-\tfrac12\bigl[(t-1)^2-(T(x^2)-1)\bigr]$
(the inner functions $T(x^2),T(x^3),T(x^4)$ being analytic at $x=\rho$ since
$\rho^2<\rho$), one has the exact cancellation
\[
\frac{\partial G}{\partial t}\Big|_{x=\rho,\,t=a}
= \rho\,Z_{S_3}\bigl(a,T(\rho^2),T(\rho^3)\bigr)-(a-1)
\;=\;0
\]
by Lemma~\ref{lem:cancel} and the defining equation of $T$. Hence the
$\sqrt{1-x/\rho}$ term of $T$ is annihilated in $A$; the singular expansion of
$A$ begins at the $(1-x/\rho)^{3/2}$ term, and singularity analysis
\cite[VI--VII]{FS2009} yields, for every $m$, the expansion \eqref{eq:asym} with
exponent $n^{-5/2}$.
\end{proposition}

\begin{proof}[Comments on the proof]
The square-root type of $T$ at $\rho$ is the standard smooth-implicit-function
schema \cite[Thm.~VII.3]{FS2009} applied to \eqref{eq:T}, the characteristic
condition being the second displayed equation (verified numerically to $55$
digits by a two-dimensional Newton iteration in $(\rho,a)$; the system is
analytic and nondegenerate there, $\det J \neq 0$). The displayed cancellation
is exact and is the celebrated mechanism of Otter \cite{Otter1948}: unrooting
shifts the polynomial growth from $n^{-3/2}$ (rooted) to $n^{-5/2}$ (unrooted).
This is the structural reason the constants of \eqref{eq:asym} exist; it is also
the reason \emph{no new} closed-form structure appears at the dominant
singularity --- the local behavior is generic square-root, indistinguishable
from that of any other simple variety of trees, with all arithmetic specificity
pushed into the constants $\rho,a,b,C,c_1,\dots$, none of which is known (or
expected) to be expressible in classical constants.
\end{proof}

\subsubsection{The cascade and the (conjectural) natural boundary}

\begin{proposition}[Singularity cascade; sketch]\label{prop:cascade}
The functional equation \eqref{eq:T} links $T(x)$ to $T(x^2)$ and $T(x^3)$.
Beyond $\rho$, the analytic continuation of $T$ supplied by \eqref{eq:T}
acquires singularities at points where the \emph{inner} evaluations become
singular: $x^2=\rho$ ($x=\pm\rho^{1/2}$, $|x|\approx0.5959$), $x^3=\rho$ (three
points, $|x|\approx0.7084$), and inductively at solutions of
$x^{2^a3^b}=\rho$, of moduli $\rho^{1/(2^a3^b)}\to1^-$, with arguments
equidistributing on the circle. Granting that no unexpected cancellation kills
any of these (each finite instance is numerically checkable; the cancellation at
$\rho$ itself, Proposition~\ref{prop:otter}, is the unique structurally forced
one), the singularities of $T$, hence of $A$, accumulate densely on $|x|=1$:
the unit circle is a natural boundary.
\end{proposition}

\begin{corollary}[Conditional global non-holonomicity]\label{cor:conditional}
If Proposition~\ref{prop:cascade} holds, then $A$ has infinitely many
singularities; since a D-finite function has only finitely many (the roots of
its leading coefficient) \cite{Stanley1980}, $a(n)$ is not P-recursive for
\emph{any} order and degree, and Corollary~\ref{cor:noWZ} holds with the region
replaced by ``everywhere''. The certificates of \S\ref{sec:holo} are exactly the
unconditional finite shadow of this statement.
\end{corollary}

\subsubsection{No transformation theory}

The substitution monoid acting in \eqref{eq:T}--\eqref{eq:A} is generated by
$x\mapsto x^2$ and $x\mapsto x^3$: a \emph{mixed Mahler} structure. For pure
$k$-Mahler equations there is a developed Galois/rigidity theory; its flagship
results --- e.g.\ a series satisfying both a $p$-Mahler and a $q$-Mahler
equation with multiplicatively independent $p,q$ must be rational
\cite{AdamczewskiBell2017,SchafkeSinger2019} --- assert precisely that the
$\{2,3\}$-mixed situation admits \emph{no} common structure except in trivial
cases. This is the antithesis of the modular situation of
Remark~\ref{rem:rademacher}: $\eta$'s boundary singularities are all
$\mathrm{SL}_2(\ZZ)$-equivalent, while the cascade points
$\rho^{1/(2^a3^b)}\zeta$ of Proposition~\ref{prop:cascade} are equivalent under
\emph{nothing}. A Rademacher-type exact formula for $a(n)$ would require an
exact local expansion at every cascade point \emph{plus} a summable global
organization of those expansions; no mathematical technology supplying either
ingredient currently exists.

%======================================================================
\section{Necessary conditions on any solution of (A)}\label{sec:necessary}
%======================================================================

\begin{theorem}[Profile of a hypothetical closed form]\label{thm:profile}
Any $F$ solving Question~\ref{q:A} (even only for all $n\ge n_0$, with
$n_0\le 25$, say) must satisfy \emph{all} of the following.
\begin{enumerate}
\item \textbf{Outside the WZ core.} If $F\in\mathcal{C}$
(Definition~\ref{def:class}), its minimal (inhomogeneous) annihilating operator
must have $(r,d)$ \emph{outside} the entire staircase of Table~\ref{tab:holo}
(after the degree shift $d\mapsto d+n_0-1$ of Lemma~\ref{lem:window}); in
particular order $\ge3$ even at degree $150$, order $\ge305$ at constant
coefficients. No formula of comparable complexity has ever occurred in
enumerative combinatorics; conditionally on Proposition~\ref{prop:cascade},
$F\notin\mathcal{C}$ outright.
\item \textbf{Not algebraic, not Abelian.} The generating function of $F$ can
satisfy no $P(x,y)=0$ within Table~\ref{tab:alg} and no first-order ADE of
degree $\le7$ per variable.
\item \textbf{Not a digit formula.} $F \bmod 2$ requires $\ge31$ states of
$2$-kernel, $F\bmod3$ requires $\ge40$ of $3$-kernel; $F\bmod m$ is aperiodic
(period $>200$) for $m\le7$. So $F$ cannot factor through any small
finite-state functional of the digits of $n$.
\item \textbf{Must implement the Burnside collapse.} $F$ must reproduce
\eqref{eq:burnside}, a partition-indexed sum of unbounded dimension, within a
bounded-depth expression --- the only known precedent for which (Rademacher)
consumed a modular transformation group that the alkane system provably lacks in
its pure-Mahler shadow and is not known to possess in any form
(\S\ref{sec:structure}).
\item \textbf{Exponential precision.} By Proposition~\ref{prop:rounding}, $F$
must beat the truncated asymptotic scale \eqref{eq:asym} by a factor
$\alpha^{-n}n^{5/2}$: it must encode the constants
$\rho, C, c_1, c_2,\dots$ to unbounded accuracy, i.e.\ contain the full
transcendental content of the singular expansion, not any finite part of it.
\end{enumerate}
\end{theorem}

\begin{remark}[What remains genuinely open]
Nothing above excludes: (i) exotic closed forms outside $\mathcal{C}$ (e.g.\
involving $\lfloor\cdot\rfloor$ of irrational multiples, or special-function
evaluations with $n$ in analytically nontrivial positions) --- though
items (3) and (5) constrain these severely; (ii) a future
``mixed-Mahler circle method'' supplying the missing transformation theory;
(iii) larger certified perimeters (the method of \S\ref{sec:method} scales to
thousands of terms with FFT-based convolution and block elimination). We regard
(ii) as the genuine mathematical content of Question~\ref{q:A}: an exact formula
for the alkane numbers is equivalent, in practice, to an exact asymptotic theory
for mixed Mahler systems --- a theory whose first nontrivial instance would
itself be a major result.
\end{remark}

%======================================================================
\appendix
\section{Reproducibility}\label{app:repro}
%======================================================================

All computations were performed in Python (exact integer arithmetic; modular
linear algebra over $\FF_{p_1},\FF_{p_2}$ with $64$-bit safe word size; mpmath
at $60$ decimal digits for the analytic constants only, which feed no
certificate). Scripts (in the \texttt{alkane\_asym} project directory):

\begin{description}
\item[\texttt{asym.py}] Computes $t(0..620)$, $a(0..620)$ from
\eqref{eq:T}--\eqref{eq:A} with exactness assertions; Newton iteration for
$(\rho,a)$ to $55$ digits; Richardson--Neville extrapolation for
$C,c_1,c_2,c_3$ on two disjoint node sets (agreement: $C$ to $1.7\times10^{-18}$,
$c_1$ to $10^{-11}$, $c_2$ to $2\times10^{-11}$).
\item[\texttt{solveA.py}] The certificate battery of
\S\S\ref{sec:holo}--\ref{sec:autom} and the labeled-tree verification of
\S\ref{sec:labeled}, including all decoys (Catalan recurrence and algebraic
equation; Thue--Morse kernel) and the second-prime re-checks.
\item[\texttt{verify\_bfile.py}] Term-by-term comparison of $a(0..620)$ against
the OEIS \texttt{b000602.txt} ($1001$ independent terms): \emph{all $621$ values
agree}.
\item[\texttt{crosscheck.py}] Reproduction of Kotesovec's published constant for
\seqT{} to $22$ digits by the same pipeline.
\end{description}

Indicative cost: series to $n=620$ in $O(N^2)$ word-multiplications of big
integers (seconds); each rank certificate is one dense elimination of a
$\le650\times\le610$ matrix over $\FF_p$ (milliseconds); the full battery runs
in well under a minute on a laptop.

%======================================================================
\begin{thebibliography}{99}
%======================================================================

\bibitem{AdamczewskiBell2017}
B.~Adamczewski and J.~P. Bell,
\emph{A problem about Mahler functions},
Ann. Sc. Norm. Super. Pisa Cl. Sci. \textbf{17} (2017), 1301--1355.

\bibitem{AlloucheShallit2003}
J.-P. Allouche and J.~Shallit,
\emph{Automatic Sequences: Theory, Applications, Generalizations},
Cambridge University Press, 2003.

\bibitem{Cayley1875}
A.~Cayley,
\emph{\"Uber die analytischen Figuren, welche in der Mathematik B\"aume genannt
werden und ihre Anwendung auf die Theorie chemischer Verbindungen},
Ber. Dtsch. Chem. Ges. \textbf{8} (1875), 1056--1059.

\bibitem{Cobham1969}
A.~Cobham,
\emph{On the base-dependence of sets of numbers recognizable by finite
automata}, Math. Systems Theory \textbf{3} (1969), 186--192.

\bibitem{Comtet1964}
L.~Comtet,
\emph{Calcul pratique des coefficients de Taylor d'une fonction alg\'ebrique},
Enseign. Math. \textbf{10} (1964), 267--270.

\bibitem{Finch2003}
S.~R. Finch,
\emph{Mathematical Constants},
Cambridge University Press, 2003, \S5.6.1 (Otter's tree enumeration constants /
chemical isomers).

\bibitem{FS2009}
P.~Flajolet and R.~Sedgewick,
\emph{Analytic Combinatorics},
Cambridge University Press, 2009 (alkanes: p.~477--478; transfer theorems:
Ch.~VI--VII).

\bibitem{HararyNorman1960}
F.~Harary and R.~Z. Norman,
\emph{Dissimilarity characteristic theorems for graphs},
Proc. Amer. Math. Soc. \textbf{11} (1960), 332--334.

\bibitem{Lipshitz1989}
L.~Lipshitz,
\emph{D-finite power series},
J. Algebra \textbf{122} (1989), 353--373.

\bibitem{Moon1970}
J.~W. Moon,
\emph{Counting Labelled Trees},
Canadian Mathematical Monographs, 1970.

\bibitem{OEIS}
OEIS Foundation Inc.,
\emph{The On-Line Encyclopedia of Integer Sequences},
sequences A000602, A000598; $b$-file \texttt{b000602.txt} (P.~von Br\"omssen,
terms $0..1000$). \url{https://oeis.org}

\bibitem{Otter1948}
R.~Otter,
\emph{The number of trees},
Ann. of Math. \textbf{49} (1948), 583--599.

\bibitem{PWZ1996}
M.~Petkov\v{s}ek, H.~S. Wilf and D.~Zeilberger,
\emph{A=B},
A K Peters, 1996.

\bibitem{Polya1937}
G.~P\'olya,
\emph{Kombinatorische Anzahlbestimmungen f\"ur Gruppen, Graphen und chemische
Verbindungen},
Acta Math. \textbf{68} (1937), 145--254.

\bibitem{Rademacher1937}
H.~Rademacher,
\emph{On the partition function $p(n)$},
Proc. London Math. Soc. \textbf{43} (1937), 241--254.

\bibitem{RainsSloane1999}
E.~M. Rains and N.~J.~A. Sloane,
\emph{On Cayley's enumeration of alkanes (or 4-valent trees)},
J. Integer Seq. \textbf{2} (1999), Article 99.1.1.

\bibitem{RHB1976}
R.~W. Robinson, F.~Harary and A.~T. Balaban,
\emph{The numbers of chiral and achiral alkanes and monosubstituted alkanes},
Tetrahedron \textbf{32} (1976), 355--361.

\bibitem{SchafkeSinger2019}
R.~Sch\"afke and M.~F. Singer,
\emph{Consistent systems of linear differential and difference equations},
J. Eur. Math. Soc. \textbf{21} (2019), 2751--2792.

\bibitem{Stanley1980}
R.~P. Stanley,
\emph{Differentiably finite power series},
European J. Combin. \textbf{1} (1980), 175--188.

\bibitem{WilfZeilberger1992}
H.~S. Wilf and D.~Zeilberger,
\emph{An algorithmic proof theory for hypergeometric (ordinary and ``$q$'')
multisum/integral identities},
Invent. Math. \textbf{108} (1992), 575--633.

\end{thebibliography}

\end{document}
